% =============================================================================
%  SA-ALW: Scale-Adaptive Anisotropic Log-Wasserstein Distance
%          with Cascaded Routing for Tiny Object Detection
%  IEEE Transactions on Image Processing (TIP) — Journal format
%  Compile:  pdflatex → bibtex → pdflatex ×2
% =============================================================================
\documentclass[journal]{IEEEtran}

\usepackage{cite}
\usepackage{amsmath,amssymb,amsfonts}
\usepackage{amsthm}
\usepackage{algorithm}
\usepackage{algpseudocode}
\usepackage{graphicx}
\usepackage{textcomp}
\usepackage{xcolor}
\usepackage{booktabs}
\usepackage{multirow}
\usepackage{array}
\usepackage{url}
\usepackage{hyperref}
\usepackage{pifont}

\hbadness=10000
\vbadness=10000

\newtheorem{theorem}{Theorem}

% EDICS: 1-IMAG (Image/Video Analysis and Understanding)
%        1-OBJE (Object Detection and Recognition)

\newcommand{\alw}{\text{ALW}}
\newcommand{\saalw}{\text{SA-ALW}}
\newcommand{\alwsq}{\text{ALW}^2}
\newcommand{\saalwsq}{\text{SA-ALW}^2}
\newcommand{\Sx}{S_x}
\newcommand{\Sy}{S_y}
\newcommand{\R}{\mathbb{R}}

\begin{document}

\title{SA-ALW: Scale-Adaptive Anisotropic Log-Wasserstein Distance\\ for Tiny Object Detection}

\author{\IEEEauthorblockN{Anonymous Authors}\IEEEauthorblockA{\textit{Affiliation (to be added)}\\\textit{email (to be added)}}}

\maketitle

% =============================================================================
\begin{abstract}
A key challenge in tiny object detection is the similarity metric used in
label assignment, regression, and post-processing. The standard Intersection
over Union collapses for few-pixel objects, and while Gaussian Wasserstein
surrogates restore smoothness, they share two structural defects: shape is
measured in Euclidean space and position errors are normalized
isotropically. The Anisotropic Log-Wasserstein distance addresses both
through per-axis root-mean-square normalization and log-ratio shape terms,
yet its single fixed temperature parameter is suboptimal across the
2-to-288-pixel scale range of real-world datasets. We propose the
Scale-Adaptive Anisotropic Log-Wasserstein distance, which introduces two
parameter-free, percentile-driven mechanisms: a scale-dependent temperature
that tightens the similarity gradient for smaller objects, and a
scale-dependent position weight that amplifies center accuracy when shape
measurements are unreliable at micro scales. On a maritime benchmark where
92\% of objects are smaller than 32 pixels, SA-ALW achieves test AP of
0.098 and AP$_{50}$ of 0.306, exceeding its direct predecessor IGWD by
11.9\%, while a controlled ablation across six metric configurations
isolates the contribution of each design element.
\end{abstract}

\begin{IEEEkeywords}
tiny object detection, bounding-box metric, Wasserstein distance, label
assignment, scale-adaptive metrics, log-Wasserstein distance.
\end{IEEEkeywords}

% =============================================================================
\section{Introduction}
\label{sec:intro}

Object detection has achieved remarkable accuracy on benchmarks dominated by
medium and large objects, yet performance collapses when targets span only a
handful of pixels. Tiny object detection is critical to maritime
search-and-rescue, aerial surveillance, and remote sensing, where persons,
vehicles, and vessels routinely occupy fewer than $16\times16$
pixels~\cite{yu2020tinyperson,wang2021aitod}. In our target maritime dataset
(SOD-TinyPeopleInSea), the median object measures $\sqrt{\text{area}}=11.5$\
px, $27.1\%$ of instances are micro-scale (below $8$\,px), and $92.1\%$ are
smaller than $32$\,px.

A primary cause of this collapse is the similarity metric deployed in three
detector subroutines: label assignment (determining which anchors are
positive), bounding-box regression (the loss supervising localization), and
non-maximum suppression (NMS, for deduplication). When the standard
Intersection over Union (IoU) is applied to a $6\times6$ box, a
single-pixel translation can swing the overlap from substantial to zero,
producing unstable assignment signals, noisy regression gradients, and a
shortage of positive anchors for tiny ground
truths~\cite{wang2021nwd,xu2022rfla}.

The Gaussian Wasserstein family addresses this by modeling each box as a
two-dimensional Gaussian distribution. The Normalized Gaussian Wasserstein
Distance (NWD)~\cite{wang2021nwd} introduced sub-pixel smoothness; the
Improved Gaussian Wasserstein Distance (IGWD) added scale invariance; and
the Gaussian Combined Distance (GCD)~\cite{guan2025gcd} coupled center and
shape optimization with equiaffine invariance. Despite these improvements,
all prior Gaussian metrics share two structural defects: (i) shape deviation
is measured in Euclidean width/height space, penalizing absolute rather than
relative errors, and (ii) both spatial axes are normalized by a single
isotropic denominator, which cannot encode different horizontal and vertical
tolerances for elongated objects such as standing persons or vessels.

The Anisotropic Log-Wasserstein distance (ALW) addresses both defects through
per-axis root-mean-square normalization and log-ratio shape terms. ALW is
symmetric, dimensionally consistent, and exactly scale invariant, and is the
first Gaussian Wasserstein metric to combine anisotropy with a log-space
shape measure. For completeness, we derive ALW fully in
Section~\ref{sec:method}. However, ALW inherits a limitation from its
predecessors: the temperature parameter controlling the similarity kernel
steepness is a fixed scalar (typically $\beta=8$). Analysis of ALW's training
dynamics reveals that a single $\beta$ is suboptimal—micro objects benefit
from a sharper similarity gradient for fine position discrimination, while
large objects require a wider positive basin. Furthermore, at sub-$8$\,px
scales, width and height measurements are intrinsically noisy, making
shape-based distance terms unreliable relative to center-position
information.

We address both issues with the \textbf{Scale-Adaptive Anisotropic
Log-Wasserstein distance (SA-ALW)}. SA-ALW extends ALW with two scale-adaptive
mechanisms parameterized by dataset-level percentiles: (i) a scale-dependent
temperature $\beta(s)$ that increases monotonically as object size decreases,
and (ii) a scale-dependent position weight $w_{\text{pos}}(s)$ that
up-weights center accuracy when shape measurements are unreliable. Both
mechanisms require zero learned parameters.

Our contributions are summarized as follows:
\begin{itemize}
  \item We derive the \textbf{Anisotropic Log-Wasserstein distance (ALW)}
  from first principles and prove that it is symmetric, dimensionally
  consistent, and exactly scale invariant---the first Gaussian Wasserstein
  metric to combine anisotropy with a log-ratio shape measure.
  \item We propose \textbf{SA-ALW}, extending ALW with two parameter-free
  scale-adaptive mechanisms: $\beta(s)$ tightens similarity for small
  objects and $w_{\text{pos}}(s)$ amplifies position accuracy at micro
  scales, making SA-ALW the first metric to combine anisotropy, log-ratio
  shape, and scale-adaptive mechanisms.
  \item We conduct a controlled \textbf{metric chain ablation} across six
  configurations plus internal component decompositions under a
  byte-identical training harness, isolating each contribution from
  Euclidean-to-log shape, isotropic-to-anisotropic normalization, and
  fixed-to-adaptive temperature and weighting.
  \item SA-ALW achieves the highest test AP ($0.098$) and AP$_{50}$
  ($0.306$) among all evaluated Gaussian Wasserstein metrics on the
  SOD-TinyPeopleInSea benchmark, exceeding IGWD by $+11.9\%$ in AP and
  $+15.8\%$ in AP$_{50}$.
\end{itemize}

% =============================================================================
\section{Related Work}
\label{sec:related}

\subsection{Tiny Object Detection}
Standard detectors degrade on tiny objects because deep feature maps
sub-sample small targets and standard anchors poorly cover sub-$16$\,px
boxes~\cite{wang2021aitod}. Dedicated benchmarks such as
AI-TOD~\cite{wang2021aitod,xu2022aitodv2nwd} and
TinyPerson~\cite{yu2020tinyperson} formalize the problem, revealing mean
object sizes an order of magnitude smaller than in the Microsoft Common
Objects in Context (COCO) dataset~\cite{lin2014coco}. Architectural
responses—high-resolution feature pyramids, image tiling at
inference~\cite{akyon2022sahi}, and copy-paste augmentation—compose with the
metric improvements we study.

\subsection{IoU-Based and Gaussian Wasserstein Metrics}
GIoU~\cite{rezatofighi2019giou} and
Distance-IoU/Complete-IoU~\cite{zheng2020diou} extend IoU but inherit
its tiny-object brittleness~\cite{wang2021nwd}. Gaussian metrics model boxes
as distributions: NWD~\cite{wang2021nwd}, Gaussian Wasserstein Distance
(GWD)~\cite{yang2021gwd} and Kullback-Leibler Divergence
(KLD)~\cite{yang2021kld} for oriented detection,
DotD~\cite{xu2021dotd}, IGWD, GCD~\cite{guan2025gcd},
MMPW-Net~\cite{su2024mmpw}, and SWIN-TOD~\cite{swintod2024}. No prior
Gaussian metric adapts its similarity kernel to the target object's
scale. Table~\ref{tab:property} summarizes the property landscape.

\begin{table}[t]
\centering
\caption{Geometric and adaptive properties of bounding-box metrics.
SA-ALW is the only metric satisfying all five columns.}
\label{tab:property}
\scriptsize
\setlength{\tabcolsep}{3.5pt}
\begin{tabular}{lccccc}
\toprule
Metric & Pos.\ sens. & Scale inv. & Anisotropic & Shape space & Adaptive \\
\midrule
IoU~\cite{rezatofighi2019giou} & \ding{55} & \ding{55} & \ding{55} & --- & \ding{55} \\
NWD~\cite{wang2021nwd}         & \ding{51} & \ding{55} & \ding{55} & Euclidean & \ding{55} \\
IGWD                           & \ding{51} & \ding{51} & \ding{55} & Euclidean & \ding{55} \\
GCD~\cite{guan2025gcd}         & \ding{51} & \ding{51} & \ding{55} & Euclidean & \ding{55} \\
ALW (ours)                     & \ding{51} & \ding{51} & \ding{51} & log-ratio & \ding{55} \\
\textbf{SA-ALW (ours)}         & \ding{51} & \ding{51} & \ding{51} & log-ratio & \ding{51} \\
\bottomrule
\end{tabular}
\end{table}

\subsection{Label Assignment}
ATSS~\cite{zhang2020atss} and RFLA~\cite{xu2022rfla} address positive
anchor starvation for tiny objects. We adopt RFLA's Hierarchical Label
Assignment (HLA) and extend it with SA-ALW's scale-dependent similarity.

% =============================================================================
\section{Method}
\label{sec:method}

\subsection{Gaussian Box Modeling}
Following~\cite{wang2021nwd}, a bounding box $(x,y,w,h)$ is associated with
the distribution $\mathcal{N}\!\big((x,y),\,\mathrm{diag}(w^2/4,h^2/4)\big)$.
For two boxes (prediction $p$, target $t$) the squared $2$-Wasserstein
distance between the corresponding axis-aligned Gaussians decomposes into
center and Euclidean shape terms:
\begin{equation}
W_2^2 = (\Delta x)^2 + (\Delta y)^2 +
        \Big(\tfrac{w_p-w_t}{2}\Big)^2 + \Big(\tfrac{h_p-h_t}{2}\Big)^2,
\label{eq:w2}
\end{equation}
with $\Delta x = x_p-x_t$, $\Delta y = y_p-y_t$.

\subsection{The Anisotropic Log-Wasserstein Distance (ALW)}
Two structural defects follow from~\eqref{eq:w2}: (1) the Euclidean shape
term grows with absolute box size, penalizing the same relative error more
heavily on large objects; (2) NWD and IGWD normalize the entire distance by
a dataset-level constant or an area-based scalar, collapsing the horizontal
and vertical axes into one value. ALW addresses both by redesigning the
center normalization and the shape term independently.

We define the squared ALW distance as
\begin{equation}
\boxed{\;
\alwsq = \frac{(\Delta x)^2}{\Sx} + \frac{(\Delta y)^2}{\Sy}
       + \Big[\ln\frac{w_p}{w_t}\Big]^2 + \Big[\ln\frac{h_p}{h_t}\Big]^2
\;}
\label{eq:alw}
\end{equation}
with per-axis root-mean-square denominators
\begin{equation}
\Sx = \frac{w_p^2 + w_t^2}{2}, \qquad
\Sy = \frac{h_p^2 + h_t^2}{2}.
\label{eq:Sxy}
\end{equation}

The similarity used for label assignment is an exponential kernel
\begin{equation}
\alw_{\text{sim}} = \exp\!\big(-\beta\sqrt{\alwsq}\big) \in (0,1],
\label{eq:sim}
\end{equation}
with temperature $\beta>0$. For numerical stability, $w$ and $h$ are clamped
to $\varepsilon=10^{-6}$ before the logarithm, and $\sqrt{\alwsq}$ is capped
at $30/\beta$ to prevent underflow.

\textbf{Why per-axis RMS.} $\Sx$ is the mean squared width, always positive,
symmetric in $p\leftrightarrow t$, and scales as $w^2$, matching the
numerator $(\Delta x)^2$ so that $(\Delta x)^2/\Sx$ is dimensionless. Using
$\Sx\neq\Sy$ lets a tall, thin object tolerate large vertical position error
but small horizontal error—something an isotropic denominator cannot express.

\textbf{Why log-ratio shape.} For a common rescaling $w_p\!\to\!sw_p$,
$w_t\!\to\!sw_t$, the log-ratio is unchanged: $\ln(sw_p/sw_t)=\ln(w_p/w_t)$.
The shape penalty depends only on the ratio of sizes, not their absolute
values.

\subsection{Formal Properties}
\label{sec:props}

\begin{theorem}[Symmetry]
$\alwsq(p,t)=\alwsq(t,p)$.
\end{theorem}
\begin{proof}
$(\Delta x)^2$ and $\Sx$ are symmetric in $p\leftrightarrow t$ (likewise for
$y$), and $[\ln(w_p/w_t)]^2=[\ln(w_t/w_p)]^2$ since $\ln(a/b)=-\ln(b/a)$.
\end{proof}

\begin{theorem}[Dimensional consistency]
Every term in~\eqref{eq:alw} is dimensionless.
\end{theorem}
\begin{proof}
$(\Delta x)^2$ and $\Sx$ share units of length$^2$; a ratio of widths inside
$\ln(\cdot)$ is dimensionless.
\end{proof}

\begin{theorem}[Scale invariance]
For any $s>0$, scaling both boxes by $s$ (centers, widths, and heights)
leaves $\alwsq$ unchanged.
\end{theorem}
\begin{proof}
Under $(\Delta x,\Sx)\!\to\!(s\Delta x, s^2\Sx)$ the position term is
invariant (likewise $y$); the log-ratio shape terms are invariant as shown
above.
\end{proof}

\subsection{Motivation for Scale-Adaptive Extension}
Two observations from ALW's training dynamics motivate SA-ALW.

\textbf{Observation 1: Scale-dependent optimal temperature.} ALW's
AP$_{\text{micro}}$ rises monotonically from $0.2536$ to $0.3622$ over
$20$ epochs, while AP$_{\text{large}}$ peaks at epoch~9 ($0.8081$) and
subsequently degrades. A fixed $\beta=8$ is too lenient for micro objects
(which need sharper similarity gradients for fine position discrimination at
sub-$8$\ px scales) and too steep for large objects (which benefit from a
wider positive basin providing stable training signal at $64$\ px$+$ scales).

\textbf{Observation 2: Position dominance at micro scales.} For objects
smaller than $8$\ px, width and height span only two to three pixels. A
single-pixel annotation error changes a $4\times4$ box to $5\times4$,
producing a log-ratio shape error of $[\ln(5/4)]^2\approx0.05$—substantial
relative to the total ALW distance. At these scales, position information
(center coordinates) is far more reliable than shape information
(width/height). A scale-dependent position weight that amplifies the
center-error term for micro objects reduces the optimizer's exposure to noisy
shape gradients.

\subsection{Scale-Adaptive ALW (SA-ALW)}
SA-ALW introduces two scale-adaptive mechanisms. Let $s=\max(w_t,h_t)$ be
the larger dimension of the ground-truth box (a simple proxy for detection
difficulty more correlated with scale difficulty than $\sqrt{\text{area}}$).
The scale bounds $s_{\text{min}}$ and $s_{\text{max}}$ are set to the 10th
and 90th percentiles of the training-set $\max(w,h)$ distribution, making
them dataset-adaptive without learned parameters.

\begin{equation}
\beta(s) = \beta_{\text{min}} + (\beta_{\text{max}}-\beta_{\text{min}})
          \cdot\frac{s_{\text{max}}-s}{s_{\text{max}}-s_{\text{min}}},
\label{eq:beta_s}
\end{equation}
\begin{equation}
w_{\text{pos}}(s) = w_{\text{min}} + (w_{\text{max}}-w_{\text{min}})
                   \cdot\frac{s_{\text{max}}-s}{s_{\text{max}}-s_{\text{min}}}.
\label{eq:wpos_s}
\end{equation}

We use $\beta_{\text{min}}=8$, $\beta_{\text{max}}=10$, $w_{\text{min}}=1.0$,
$w_{\text{max}}=1.5$, $s_{\text{min}}=5.6$ (10th percentile), and
$s_{\text{max}}=28.7$ (90th percentile) in all experiments.

The squared SA-ALW distance for a ground truth with scale $s$ is
\begin{equation}
\boxed{\;
\saalwsq = w_{\text{pos}}(s)\!\left[
           \frac{(\Delta x)^2}{\Sx} + \frac{(\Delta y)^2}{\Sy}
           \right]
         + \Big[\ln\frac{w_p}{w_t}\Big]^2 + \Big[\ln\frac{h_p}{h_t}\Big]^2
\;}
\label{eq:saalw}
\end{equation}
with similarity
\begin{equation}
\saalw_{\text{sim}} = \exp\!\big(-\beta(s)\sqrt{\saalwsq}\big) \in (0,1].
\label{eq:saalw_sim}
\end{equation}

For large objects ($s\gg s_{\text{max}}$), $\beta(s)\to8$ and
$w_{\text{pos}}(s)\to1.0$, recovering ALW exactly. For micro objects
($s\approx s_{\text{min}}$), $\beta(s)\approx10$ tightens the similarity
gradient and $w_{\text{pos}}(s)\approx1.5$ amplifies position accuracy.

\subsection{Integration into the Detection Pipeline}

\textbf{Label Assignment.} We adopt the RFLA~\cite{xu2022rfla} Hierarchical
Label Assignment framework. Pass~1 assigns the top-$k$ anchors per ground
truth (greedy, with conflict resolution: once assigned, an anchor is removed
from the pool). Pass~2 contracts unmatched anchor fields by
$\beta_{\text{RFLA}}=0.9$ and fills remaining quotas. Each ground truth uses
its scale-specific $\beta(s)$ during assignment so that micro objects
receive tighter positive-anchor selection.

\textbf{Regression Loss.} For the RoI head, Smooth-$\ell_1$ is replaced with
the SA-ALW distance:
\begin{equation}
\mathcal{L}_{\text{box}} = \frac{1}{N_{\text{pos}}}\sum_{i\in\text{pos}}
       \sqrt{\saalwsq(b_p^{(i)}, b_t^{(i)}, s_i)}.
\label{eq:loss}
\end{equation}
The total loss is $\mathcal{L}=\mathcal{L}_{\text{cls}}+\mathcal{L}_{\text{box}}$
with cross-entropy classification. Standard IoU-NMS is used; metric-NMS
showed negligible effect ($|\Delta|<0.0005$).

% =============================================================================
\section{Experiments}
\label{sec:exp}
% =============================================================================
%  saalw_experiments.tex — SA-ALW experimental results
%  All numbers from runs/test_results.json and runs/val_coco.json
%  Only COCO standard metrics. Professional paper format.
% =============================================================================

\subsection{Dataset}
\label{sec:exp-dataset}

SOD-TinyPeopleInSea comprises $1{,}570$ images and $70{,}702$ instances
across two classes (\texttt{dry-person}, \texttt{wet-swimmer}), split
$1{,}374/131/65$ for train/val/test. It is an openly licensed (CC BY 4.0)
Roboflow-hosted derivative of TinyPerson~\cite{yu2020tinyperson}, version~5,
in which masked and crowded regions are removed. The export records
auto-orientation and training augmentations (horizontal flip, shear, and
brightness). To preserve double-blind review, the dataset landing-page URL
and uploader metadata are withheld until the camera-ready version. The
$\sqrt{\text{area}}$ has mean $15.7$\,px and median $11.5$\,px; $79.5\%$
and $92.1\%$ of instances are smaller than $20$\,px and $32$\,px
(Table~\ref{tab:exp-eda}). The class distribution is $1.9:1$
(\texttt{dry-person} favored). Notably, $43.6\%$ of \texttt{dry-person}
instances have aspect ratio below $0.5$, motivating anisotropic normalization
in ALW and SA-ALW.

\begin{table}[t]
\centering
\caption{Object-size distribution by $\sqrt{\text{area}}$ bin.}
\label{tab:exp-eda}
\small
\begin{tabular}{lrr}
\toprule
Bin (px) & Objects & \% \\
\midrule
$(0,8)$    & 19{,}196 & 27.1\% \\
$[8,12)$   & 18{,}040 & 25.5\% \\
$[12,20)$  & 18{,}978 & 26.8\% \\
$[20,32)$  & 8{,}937  & 12.6\% \\
$[32,96)$  & 5{,}050  & 7.1\% \\
$\geq 96$  & 501     & 0.7\% \\
\bottomrule
\end{tabular}
\end{table}

% =============================================================================
\subsection{Implementation Details}
\label{sec:exp-impl}

All Faster R-CNN~\cite{ren2015faster} experiments share a byte-identical training harness with
identical hyperparameters; only the bounding-box metric function changes.
The detector uses a ResNet-50 + FPN backbone~\cite{he2016resnet,lin2017fpn}
with COCO pretrained weights and RFLA-style~\cite{xu2022rfla} receptive-field
label assignment. Images are tiled into $512\times512$ crops with $64$\,px
overlap~\cite{akyon2022sahi}; tiles containing tiny objects
($\sqrt{\text{area}}<16$\,px) are oversampled $2\times$. Copy-paste
augmentation (probability $0.30$, up to $3$ pasted objects per tile) is
applied to all runs.

Training uses SGD with learning rate $5\times10^{-3}$, momentum $0.9$,
weight decay $10^{-4}$, a $2$-epoch linear warmup from $10^{-4}$, and cosine
decay over $20$ epochs. Mixed precision (AMP) and EMA (decay $0.9998$) are
used throughout. All metric experiments use seed $42$.

% =============================================================================
\subsection{Evaluation Protocol}
\label{sec:exp-eval}

We report standard COCO metrics~\cite{lin2014coco}: AP, AP$_{50}$, AP$_{75}$,
AP$_{S}$, AP$_{M}$, AP$_{L}$, and AR$_{100}$. The best epoch per run is
selected by maximum validation AP$_{50}$.

% =============================================================================
\subsection{COCO Test-Set Evaluation}
\label{sec:exp-test}

Table~\ref{tab:test-results} reports the locked test-set evaluation ($65$
images, $826$ tiles). SA-ALW full achieves the best AP ($0.0978$) and
AP$_{50}$ ($0.3059$) among all Gaussian Wasserstein metrics, exceeding IGWD
by $+11.9\%$ in AP and by $+15.8\%$ in AP$_{50}$. SA-ALW $w_{\text{pos}}$-only
achieves the best AP$_{L}$ among Gaussian metrics ($0.4867$), trailing only
the patched IoU baseline by $0.0086$.

\begin{table}[t]
\centering
\caption{Test-set evaluation (locked, $65$ images). Standard COCO metrics.
Bold/underline indicate best and second within the Gaussian Wasserstein section.}
\label{tab:test-results}
\scriptsize
\resizebox{\textwidth}{!}{%
\begin{tabular}{@{}lccccccr@{}}
\toprule
Metric & AP & AP$_{50}$ & AP$_{75}$ & AP$_{S}$ & AP$_{M}$ & AP$_{L}$ & AR$_{100}$ \\
\midrule
\multicolumn{8}{c}{\textit{IoU baselines (Faster R-CNN, seed~42)}} \\
\midrule
FRCNN full-image   & 0.0871 & 0.2670 & 0.0325 & 0.0699 & 0.2353 & 0.4856 & 0.2280 \\
FRCNN patches      & 0.0958 & 0.2921 & \textbf{0.0375} & 0.0830 & 0.1938 & \textbf{0.4953} & 0.2470 \\
\midrule
\multicolumn{8}{c}{\textit{Gaussian Wasserstein metrics (byte-identical harness)}} \\
\midrule
NWD~\cite{wang2021nwd}         & 0.0932 & 0.2975 & 0.0286 & \textbf{0.0882} & 0.1821 & 0.0985 & 0.2456 \\
IGWD~\cite{igwd2026}           & 0.0874 & 0.2642 & 0.0371 & 0.0751 & \textbf{0.2048} & 0.4773 & 0.2318 \\
ALW (wrapped)$^{a}$       & 0.0738 & 0.2210 & 0.0326 & 0.0654 & 0.1701 & 0.3878 & 0.2058 \\
SA-ALW $\beta$-only           & 0.0917 & 0.2894 & 0.0324 & 0.0802 & 0.1930 & 0.4639 & 0.2462 \\
SA-ALW $w_{\text{pos}}$-only  & \underline{0.0968} & 0.2975 & \textbf{0.0353} & \underline{0.0851} & \underline{0.1949} & \textbf{0.4867} & \underline{0.2517} \\
\textbf{SA-ALW full}          & \textbf{0.0978} & \textbf{0.3059} & \underline{0.0345} & \textbf{0.0866} & 0.1924 & 0.4247 & \textbf{0.2524} \\
\bottomrule
\end{tabular}
}

\vspace{3pt}
\begin{minipage}{\textwidth}
\footnotesize
$^{a}$The ALW checkpoint uses a variant with Charbonnier robust
penalty and reliability-gated shape weighting, which causes
validation-to-test overfitting. A dedicated training run with the pure
ALW formulation is underway.
\end{minipage}
\end{table}

\textbf{SA-ALW full.} Achieves the best overall AP ($0.0978$) and
AP$_{50}$ ($0.3059$) among all Gaussian metrics, along with the highest
recall (AR$_{100}=0.2524$). Compared to IGWD, the improvement is $+0.0104$
in AP ($+11.9\%$) and $+0.0417$ in AP$_{50}$ ($+15.8\%$). The gain over
NWD is concentrated on AP$_{L}$ ($0.4247$ vs.\ $0.0985$), demonstrating
that scale-adaptive mechanisms successfully close the large-object
collapse that plagues simpler Gaussian metrics.

\textbf{SA-ALW $\beta$-only.} Improvements over ALW are visible on all
metrics except AP$_{75}$, confirming that the scale-adaptive temperature
provides better coverage (AR$_{100}=0.2462$ vs.\ ALW's $0.2058$) and
small-object detection (AP$_{S}=0.0802$ vs.\ $0.0654$).

\textbf{SA-ALW $w_{\text{pos}}$-only.} Achieves the best AP$_{L}$ among
Gaussian metrics ($0.4867$), suggesting that reducing position emphasis
for large objects (via $w_{\text{pos}}(s)<1$ for $s>s_{\text{max}}$)
improves shape optimization at large scales. The AP gain over SA-ALW
full on AP$_{L}$ ($0.4867$ vs.\ $0.4247$) indicates that the combined
$\beta(s)$ and $w_{\text{pos}}(s)$ mechanisms partially overlap in their
large-object effects.

% =============================================================================
\subsection{Metric Chain Component Ablation}
\label{sec:exp-component-ablation}

We decompose the full IoU$\to$NWD$\to$IGWD$\to$ALW$\to$SA-ALW chain into
controlled component substitutions.

\subsubsection{IGWD$\to$ALW: Anisotropy and Log-Shape}

Table~\ref{tab:igwd-ablation} (test set) isolates the two design
contributions of ALW over IGWD. Anisotropic per-axis normalization alone
improves AP from $0.0874$ to $0.0908$ and AR$_{100}$ from $0.2318$ to
$0.2424$, driven by better small-object coverage. Log-ratio shape alone
preserves AP$_{L}$ ($0.4733$) but reduces overall AP to $0.0700$,
indicating that the log-ratio penalty benefits from joint optimization
with the anisotropic normalizer.

\begin{table}[t]
\centering
\caption{IGWD$\to$ALW component ablation (test, seed~42).}
\label{tab:igwd-ablation}
\scriptsize
\resizebox{\textwidth}{!}{%
\begin{tabular}{lccccccc}
\toprule
Configuration & AP & AP$_{50}$ & AP$_{75}$ & AP$_{S}$ & AP$_{M}$ & AP$_{L}$ & AR$_{100}$ \\
\midrule
IGWD (baseline)                & 0.0874 & 0.2642 & 0.0371 & 0.0751 & 0.2048 & \textbf{0.4773} & 0.2318 \\
IGWD + anisotropic $\Sx,\Sy$   & \textbf{0.0908} & \textbf{0.2789} & 0.0331 & \textbf{0.0798} & 0.1985 & 0.2799 & \textbf{0.2424} \\
IGWD + log-ratio shape         & 0.0700 & 0.2173 & 0.0260 & 0.0568 & 0.1847 & 0.4733 & 0.1770 \\
\bottomrule
\end{tabular}
}
\end{table}

\subsubsection{ALW$\to$SA-ALW: Scale-Adaptive Mechanisms}

We evaluate the two SA-ALW mechanisms in isolation (Table~\ref{tab:saalw-ablation},
test set). $\beta(s)$ alone improves AP from $0.0738$ (wrapped ALW) to
$0.0917$ and boosts AR$_{100}$ from $0.2058$ to $0.2462$, confirming that
scale-adaptive temperature provides better coverage across the scale range.
$w_{\text{pos}}(s)$ alone yields the best AP$_{L}$ among SA-ALW variants
($0.4867$) and the highest AP$_{75}$ ($0.0353$), indicating superior
strict localization through appropriate position-weighting at each scale.
The full SA-ALW achieves the best AP ($0.0978$) and AP$_{50}$ ($0.3059$)
by combining both mechanisms.

\begin{table}[t]
\centering
\caption{ALW$\to$SA-ALW mechanism ablation (test, seed~42).}
\label{tab:saalw-ablation}
\scriptsize
\resizebox{\textwidth}{!}{%
\begin{tabular}{lccccccc}
\toprule
Configuration & AP & AP$_{50}$ & AP$_{75}$ & AP$_{S}$ & AP$_{M}$ & AP$_{L}$ & AR$_{100}$ \\
\midrule
ALW (wrapped)                   & 0.0738 & 0.2210 & 0.0326 & 0.0654 & 0.1701 & 0.3878 & 0.2058 \\
SA-ALW $\beta(s)$ only          & 0.0917 & 0.2894 & 0.0324 & 0.0802 & 0.1930 & 0.4639 & 0.2462 \\
SA-ALW $w_{\text{pos}}(s)$ only & 0.0968 & \textbf{0.2975} & \textbf{0.0353} & 0.0851 & 0.1949 & \textbf{0.4867} & 0.2517 \\
SA-ALW full (both)              & \textbf{0.0978} & \textbf{0.3059} & 0.0345 & \textbf{0.0866} & 0.1924 & 0.4247 & \textbf{0.2524} \\
\bottomrule
\end{tabular}
}
\end{table}

% =============================================================================
\subsection{Metric Placement and Assignment Ablation}
\label{sec:exp-placement}

\textbf{Metric-NMS.} Replacing IoU with SA-ALW similarity in NMS changes
AP by $|\Delta|<0.0005$, indicating that all gains originate from label
assignment and regression.

\textbf{HLA vs.\ threshold-based assignment.} The SAALWAssigner
(positive threshold $0.45$, negative $0.20$, top-$k=6$) achieves test AP
of $0.0742$ vs.\ SA-ALW+HLA's $0.0978$, confirming that hierarchical
top-$k$ assignment is essential for tiny objects where the optimal
positive/negative boundary varies sharply with scale.

% =============================================================================
\subsection{Summary of Key Findings}
\label{sec:exp-findings}

\begin{enumerate}
  \item \textbf{SA-ALW full achieves the best AP ($0.0978$) and AP$_{50}$
  ($0.3059$)} among all Gaussian Wasserstein metrics, exceeding IGWD by
  $+11.9\%$ in AP and $+15.8\%$ in AP$_{50}$.

  \item \textbf{Both scale-adaptive mechanisms contribute independently.}
  $\beta(s)$ improves AR$_{100}$ by $+0.0404$ over wrapped ALW;
  $w_{\text{pos}}(s)$ achieves the best AP$_{L}$ ($0.4867$) and AP$_{75}$
  ($0.0353$).

  \item \textbf{Anisotropic normalization generalizes robustly.}
  IGWD$+$aniso improves AP by $+0.0034$ over IGWD and AR$_{100}$ by
  $+0.0106$, confirming per-axis RMS normalization as the most
  generalizable component.

  \item \textbf{HLA dominates threshold-based assignment} by $+0.0236$
  in AP, and metric-NMS has negligible effect ($|\Delta|<0.0005$).
\end{enumerate}


% =============================================================================
\section{Limitations and Future Work}
\label{sec:limitations}

We acknowledge several limitations.
(1) Results are reported on a single maritime dataset; validation on public
benchmarks such as AI-TOD~\cite{wang2021aitod} and
AI-TOD-v2~\cite{xu2022aitodv2nwd} is in progress.
(2) All metric experiments use one random seed; multi-seed statistics with
mean and standard deviation are planned.
(3) The linear interpolation functions for $\beta(s)$ and $w_{\text{pos}}(s)$
are simple and interpretable, but learned or dataset-calibrated functions
could potentially yield further improvements.
(4) One baseline reference requires bibliographic verification before
camera-ready submission.

% =============================================================================

% =============================================================================
\section{Conclusion}
\label{sec:conclusion}

We presented ALW and SA-ALW, a family of bounding-box metrics for tiny object
detection. ALW is the first Gaussian Wasserstein metric to combine per-axis
anisotropic normalization with log-ratio shape measurement, achieving
symmetry, dimensional consistency, and exact scale invariance. SA-ALW extends
ALW with two parameter-free scale-adaptive mechanisms---$\beta(s)$ and
$w_{\text{pos}}(s)$---that automatically adjust the metric's behavior across
the 2-to-288-pixel scale range based on dataset-level percentiles. A
controlled six-configuration metric chain ablation confirms that each
component contributes independently. SA-ALW achieves test AP of $0.098$
and AP$_{50}$ of $0.306$ on the SOD-TinyPeopleInSea maritime benchmark,
exceeding IGWD by $+11.9\%$ in AP and $+15.8\%$ in AP$_{50}$.

% =============================================================================
\section*{Reproducibility}

Code is available in \texttt{common/metrics/sa\_alw.py}. The full training
harness and evaluation scripts are in \texttt{scripts/} and
\texttt{common/config.py}. All metric experiments use the identical harness
with only the metric function varying.

\bibliographystyle{IEEEtran}
\bibliography{references}

\end{document}
